% Appendix A --- the Cin rays and the Matérn exponent. Added 2026-09-12 to resolve the forward
% references of v0.1 (§§1-7) into blueprint chapters 8 and 10, on the author's go: two nodes are
% transcribed verbatim, both [T] and \leanok on Lean core.
%   - lem:cin-rays (chapter 8): §2's dictionary remark names its Cin functions, and the proof of
%     lem:selfdecomposable-exponents in §7 uses its clause (2), the quantile measure attached to a
%     nonincreasing profile, and its clause (1), the growth of Cin at infinity; rem:cin-constant is
%     copied with it.
%   - prop:matern-exponent (chapter 10): §3's Matérn member rests on its clause (1) for
%     admissibility and on its clause (3) for the variance and the finiteness of every moment.
%     The remark at the end of its proof pointed at the moment criterion of chapter 10's
%     prop:moments-tails, which is not transcribed; the remark cites the criterion's source instead.
% NOT transcribed: prop:matern-density (chapter 10, [A] on A15 and A16, notready), whose statement
% references prop:moments-tails and whose transcription would bring four more cited facts into
% this paper for one sentence of §3. The density identification of the Matérn member stays a cited
% closed form in §3 (A15's anchors), and the recommended end state --- narrowing prop:two-members
% by that sentence, in the blueprint --- is recorded in notes/HANDOFF-article.md for the
% formalization session, which owns blueprint/ and Formalization/.

\section{The \texorpdfstring{$\Cin$}{Cin} rays and the Mat\'ern exponent}\label{sec:appendix}

Two computations that the body of the paper uses are collected here: the exponents of the
unit-step profiles, which the proof of Lemma~\ref{lem:selfdecomposable-exponents} uses through
the tail measure they attach to a nonincreasing profile, and the exponent of the Mat\'ern
member of \S\ref{sec:axioms}. Both belong to material outside this paper (the geometry of
the admissible cone and the families at its corners) and are stated here in the form the
body uses and in no more generality. Both are machine-checked and rest on Lean core alone.

% shared with blueprint lem:cin-rays
\begin{lemma}[the $\Cin$ rays and the superposition]\label{lem:cin-rays}
Write $\Cin(z) := \int_0^z (1 - \cos v)\,v^{-1}dv$ for $z \ge 0$, extended evenly to $\RR$.
\begin{enumerate}
\item \emph{(The rays.)} For $\tau > 0$ the profile $k = \mathbf 1_{(0,\tau)}$ is admissible,
with Gaussian coefficient $0$, and its exponent is $\omega \mapsto \Cin(\tau\omega)$. The
function $\Cin$ is even and nondecreasing on $[0,\infty)$, with $\Cin(z) \le z^2/4$ and
$\Cin(z) = \tfrac14 z^2 + O(z^4)$ as $z \to 0$; and there is a finite constant $C$ with
$\Cin(z) = \log z + C + O(1/z)$ as $z \to \infty$.
\item \emph{(The superposition.)} Every admissible $F$, with Gaussian coefficient $a$ and
profile $k$, satisfies
\begin{equation}\label{eq:cone-superposition}
  F(\omega) = a\,\omega^2 + \int_{(0,\infty)} \Cin(\tau\omega)\,\varpi(d\tau),
  \qquad \omega \in \RR,
\end{equation}
where $\varpi$ is any positive measure on $(0,\infty)$ with $\varpi((x,\infty)) = k(x)$ for
almost every $x > 0$; such a $\varpi$ exists.
\end{enumerate}
\end{lemma}

\begin{proof}
(1) Evenness of $\Cin$ needs no stipulation: the integrand $v \mapsto (1 - \cos v)/v$ is odd,
so the integral from $0$ is even, and $\Cin(-z) = \Cin(z)$. Monotonicity on $[0,\infty)$ is the
positivity of that integrand there. From $1 - \cos v \le v^2/2$ we get
$\Cin(z) \le \int_0^z v/2\,dv = z^2/4$ for $z \ge 0$, and for $z < 0$ by evenness, so the bound
holds on $\RR$. For the expansion at the origin, $|\cos v - (1 - v^2/2)| \le \tfrac{5}{96}v^4$
for $|v| \le 1$ gives $|(1-\cos v)/v - v/2| \le \tfrac{5}{96}v^3$ on $(0,1]$, and integrating
over $[0,z]$,
\[
  \Bigl|\Cin(z) - \tfrac14 z^2\Bigr| \ \le\ \tfrac{5}{96}\,z^4,
  \qquad |z| \le 1,
\]
which is the asserted $\Cin(z) = \tfrac14z^2 + O(z^4)$. For the expansion at infinity, split at
$v = 1$:
\[
  \Cin(z) = \Cin(1) + \log z - \int_1^z \frac{\cos v}{v}\,dv , \qquad z \ge 1,
\]
using $(1-\cos v)/v = v^{-1} - v^{-1}\cos v$ on $[1,z]$; and integrating by parts against the
primitive $\sin v/v$,
$\int_1^z v^{-1}\cos v\,dv = \sin z/z - \sin 1 + \int_1^z v^{-2}\sin v\,dv$. The last integrand
is dominated by $v^{-2}$, so $\int_1^\infty v^{-2}\sin v\,dv$ converges absolutely and, with
\[
  C := \Cin(1) + \sin 1 - \int_1^\infty \frac{\sin v}{v^2}\,dv ,
\]
one gets $\Cin(z) - \log z - C = -\sin z/z + \int_z^\infty v^{-2}\sin v\,dv$, whence
$|\Cin(z) - \log z - C| \le 2/z$ for $z \ge 1$. That is the expansion, with an explicit
constant in the error.

For the ray, $k = \mathbf 1_{(0,\tau)}$ is nonnegative and nonincreasing on $(0,\infty)$ and
satisfies both integrability conditions of Lemma~\ref{lem:profile-integrability} at once,
being bounded by $1$ and supported in $(0,\tau)$: $\int_0^1 x\,k(x)\,dx \le \tfrac12$ and
$\int_1^\infty k(x)\,x^{-1}dx \le \tau$. So $(0,k)$ is a pair of the form \ref{eq:sd-profile},
and its exponent is
$\int_0^\tau (1-\cos\omega x)\,x^{-1}dx = \Cin(\tau\omega)$ by the substitution $v = \omega x$,
valid for every real $\omega$ because both sides vanish at $\omega = 0$ and $\Cin$ is even.

(2) A positive measure $\varpi$ on $(0,\infty)$ with $\varpi((x,\infty)) = k(x)$ for almost
every $x > 0$ is the image of Lebesgue measure on $(0,\infty)$ under the generalised inverse
$y \mapsto \sup\{u > 0 : k(u) > y\}$, restricted to $(0,\infty)$ to discard the atom the
transform leaves at the origin when $k$ is bounded. Only two properties of $k$ are used: it is
nonincreasing, and it tends to $0$ at infinity --- and the second is forced rather than
assumed, since $k$ nonincreasing and bounded below by a positive constant would make
$\int_1^\infty k(x)\,x^{-1}dx$ diverge like the harmonic integral. What comes out is the
sandwich $k(x{+}) \le \varpi((x,\infty)) \le k(x)$, whose two sides agree off the countably
many discontinuities of $k$; that is why the tail identity is asserted almost everywhere and
not everywhere, and a right-continuous version of $k$ is neither available nor needed.

Given such a $\varpi$, \ref{eq:cone-superposition} is the layer-cake exchange
\[
  \int_{(0,\infty)} \Cin(\tau\omega)\,\varpi(d\tau)
  = \int_0^\infty \frac{1 - \cos\omega x}{x}\,\varpi((x,\infty))\,dx
  = \int_0^\infty (1 - \cos\omega x)\,k(x)\,\frac{dx}{x},
\]
in which $\Cin(\tau\omega)$ is read as $\int_0^\tau (1-\cos\omega x)x^{-1}dx$ and the exchange
is Tonelli, for which $\varpi$ is $\sigma$-finite: the tail measure of an unbounded profile
puts infinite mass near the origin, but its restriction to $[\varepsilon,\infty)$ is finite
for every $\varepsilon > 0$. Adding $a\,\omega^2$ gives
\ref{eq:cone-superposition}.
\end{proof}

% blueprint rem:cin-constant, copied
\begin{remark}[the constant]\label{rem:cin-constant}
The constant of Lemma~\ref{lem:cin-rays}(1) is Euler's constant: $\Cin(z) = \gamma_E + \log z -
\mathrm{Ci}(z)$ with $\mathrm{Ci}$ the cosine integral and $\mathrm{Ci}(z) = O(1/z)$
(\sscite{dlmf2026}, \S6.2). The identification is not used anywhere in this paper, which is
why Lemma~\ref{lem:cin-rays} asserts only that the constant is finite.
\end{remark}

\ref{eq:cone-superposition} writes every admissible exponent as a superposition of the
Gaussian ray and the unit-step exponents. That the unit-step exponents are extreme rays of
the cone does not follow from the superposition and is outside this paper; only the two
clauses above are used in it.

In clause (3) of the next proposition $B$ denotes standard Brownian motion, taken
independent of the random time $T$; the letter is the statement's, and it is not the symbol
$B = \omega F'$ of Lemma~\ref{lem:selfdecomposable-exponents}.

% shared with blueprint prop:matern-exponent
\begin{proposition}[the \Matern{} family: exponent, transforms, moments]\label{prop:matern-exponent}
For $\gamma > 0$:
\begin{enumerate}
\item \emph{(Exponent.)} $F(\omega) = \gamma\log(1 + \omega^2)$ is admissible, with Gaussian
coefficient $a = 0$ and folded profile $k(x) = 2\gamma e^{-x}$; the profile at range $\theta$ is
$k(x) = 2\gamma e^{-x/\theta}$, with exponent $\gamma\log(1 + \theta^2\omega^2)$.
\item \emph{(Transforms.)} $e^{-F(t\omega)} = (1 + t^2\omega^2)^{-\gamma}$, and the increments
have transfer functions $\bigl((1 + s^2\omega^2)/(1 + t^2\omega^2)\bigr)^{\gamma}$, rational in
$\omega$ for integer $\gamma$.
\item \emph{(Moments.)} Every moment is finite: $X_t \seq B_T$ with $T$ Gamma-distributed of
shape $\gamma$ and scale $2t^2$, so
$\mathbb{E}X_t^{2n} = (2n-1)!!\,(2t^2)^n\,\Gamma(\gamma+n)/\Gamma(\gamma)$, and in particular
$\mathbb{E}X_t^2 = 2\gamma t^2$.
\end{enumerate}
\end{proposition}

Finiteness of every moment follows from the even moments by $|x|^n \le 1 + x^{2n}$, and the
Gaussian mixture is constructed rather than cited.

\begin{proof}
(1) Fix $\theta > 0$ and write $p = 1/\theta$. The integrand
$(1 - \cos\omega x)\,e^{-px}x^{-1}$ is nonnegative on $(0,\infty)$ and, since
$1 - \cos u \le u^2/2$, is bounded there by $\tfrac12\omega^2\,x\,e^{-px}$, so it is integrable
for every $\omega$; its derivative in $\omega$ is $\sin(\omega x)\,e^{-px}$, whose modulus is at
most $e^{-px}$ \emph{uniformly in} $\omega$. Differentiation under the integral sign is
therefore legitimate at every $\omega$ with one and the same dominating function, and gives
\[
  \frac{d}{d\omega}\int_0^\infty (1 - \cos\omega x)\,e^{-px}\,\frac{dx}{x}
   = \int_0^\infty \sin(\omega x)\,e^{-px}\,dx = \frac{\omega}{p^2+\omega^2},
\]
which is also the derivative of $\tfrac12\log(1 + \omega^2/p^2)$. The two functions of $\omega$
therefore differ by a constant, and both vanish at $\omega = 0$, so
\[
  \int_0^\infty (1 - \cos\omega x)\,e^{-x/\theta}\,\frac{dx}{x}
   = \tfrac12\log\bigl(1 + \theta^2\omega^2\bigr).
\]
Hence $F(\omega) = \gamma\log(1+\theta^2\omega^2)$ is of the form \ref{eq:sd-profile} with
$a = 0$ and $k(x) = 2\gamma e^{-x/\theta}$, which is nonnegative and nonincreasing on
$(0,\infty)$. Its two integrability conditions are the two dominating functions above read on
the two halves of the range: $\int_0^1 x\,k(x)\,dx \le 2\gamma$ and
$\int_1^\infty k(x)\,x^{-1}dx \le 2\gamma\int_1^\infty e^{-x/\theta}dx < \infty$. Exhibiting the
pair $(0, 2\gamma e^{-\cdot/\theta})$ is what admissibility asks for, so the argument stops
there; the range $\theta$ is carried by the computation and needs no separate dilation step.

(2) is (1) read through the similarity form of Theorem~\ref{thm:main-characterization}: the
increment transfer function is the ratio of the endpoint transforms, the denominators being
nonvanishing.

(3) Let $\Gamma_{\gamma,2t^2}$ be the Gamma law of shape $\gamma$ and scale $2t^2$, that is the
law on $(0,\infty)$ with density $u^{\gamma-1}e^{-u/2t^2}/(\Gamma(\gamma)(2t^2)^{\gamma})$. Its
Laplace transform is $\int_0^\infty e^{-\sigma u}\,\Gamma_{\gamma,2t^2}(du) =
(1 + 2t^2\sigma)^{-\gamma}$ for $\sigma \ge 0$, one Gamma integral; so the Gaussian mixture
$\int_0^\infty g_u\,\Gamma_{\gamma,2t^2}(du)$ --- the law of $B_T$ with $T$ so distributed ---
has cosine transform $(1 + 2t^2\cdot\omega^2/2)^{-\gamma} = (1 + t^2\omega^2)^{-\gamma}$, which
by (2) is the transform of the kernel at scale $t$. Both are symmetric probability measures, so
Proposition~\ref{prop:fourier-uniqueness} identifies them and $X_t \seq B_T$. Conditioning on
$T$, $\mathbb{E}X_t^{2n} = \mathbb{E}[T^n]\,\mathbb{E}[B_1^{2n}] = (2n-1)!!\;\mathbb{E}T^n$
with $\mathbb{E}[B_v^{2n}] = (2n-1)!!\,v^n$ and
$\mathbb{E}T^n = (2t^2)^n\Gamma(\gamma+n)/\Gamma(\gamma)$; the two are the same Gamma integral
read at half-integer and at integer argument. Every moment is then finite, since
$|x|^n \le 1 + x^{2n}$, and $n = 1$ with $\Gamma(\gamma+1) = \gamma\Gamma(\gamma)$ gives
$\mathbb{E}X_t^2 = 2\gamma t^2$.

\emph{Remark.} Finiteness of every moment also follows from the moment criterion for
infinitely divisible laws --- $\mathbb{E}|X|^n < \infty$ exactly when the \Levy{} measure
integrates $|x|^n$ off a neighbourhood of the origin (\sscite{sato1999levy}, Thm.~25.3,
p.~159) --- the profile decaying exponentially. That route is machine-checked too, uses
that cited fact, and is independent of this one; it does not reach the general even moment.
\end{proof}
